\subsection{Symmetric and Exterior Powers and Algebras}

% This subsection develops symmetric and exterior powers from multilinear maps
% and tensor powers, then assembles the associated graded algebras.

\begin{definition} \label{def:special-multilinear-maps} \hlt{Symmetric, Skew-Symmetric, and Alternating Maps}\\
Let $f:\mathcal{V}^{\times n}\to\mathcal{W}$ be multilinear.
\begin{enumerate}[label=(\roman*)]
\setlength{\itemsep}{0pt}
\item $f$ is \hlt{symmetric} if interchanging any two coordinate positions does not change its value.
\item $f$ is \hlt{skew-symmetric} if interchanging any two coordinate positions multiplies its value by $-1$.
\item $f$ is \hlt{alternating} if $f(\mathbf{v}_1,\ldots,\mathbf{v}_n)=0$ whenever two coordinate entries are equal.
\end{enumerate}
If $\operatorname{char}(\mathbb{F})\neq2$, alternating and skew-symmetric are equivalent.\\
If $\operatorname{char}(\mathbb{F})=2$, skew-symmetric is the same as symmetric, and alternating is the stronger condition.
\end{definition}

\begin{definition} \label{def:permutation-action-on-tensor-powers}\label{def:symmetric-and-exterior-tensor-spaces} \hlt{Permutation Action and Symmetric/Exterior Tensor Spaces}\\
Let $\dim\mathcal{V}<\infty$, $p\in\mathbb{Z}_{\geq0}$, $\mathcal{V}^{\otimes0}=\mathbb{F}$, $T_p\defeq\mathcal{V}^{\otimes p}$, and $G_p\defeq\mathfrak{S}_p$.
\begin{enumerate}[label=(\roman*)]
\setlength{\itemsep}{0pt}
\item For $\sigma\in G_p$, the \hlt{permutation action} is $\lambda_{\sigma}:T_p\to T_p$, determined by $\lambda_{\sigma}(\mathbf{v}_1\otimes\cdots\otimes\mathbf{v}_p)=\mathbf{v}_{\sigma(1)}\otimes\cdots\otimes\mathbf{v}_{\sigma(p)}$.
\item The \hlt{symmetric tensors} are $\operatorname{ST}^{p}(\mathcal{V})\defeq\{\mathbf{z}\in T_p:\forall\sigma\in G_p,\lambda_{\sigma}\mathbf{z}=\mathbf{z}\}$.
\item The \hlt{antisymmetric tensors} are $\operatorname{AT}^{p}(\mathcal{V})\defeq\{\mathbf{z}\in T_p:\forall\sigma\in G_p,\lambda_{\sigma}\mathbf{z}=\sign(\sigma)\mathbf{z}\}$.
\item Set $R_{\vee}\defeq\operatorname{span}\{\lambda_{\sigma}\mathbf{z}-\mathbf{z}:\sigma\in G_p,\mathbf{z}\in T_p\}$.
\item The \hlt{symmetric power} is $\operatorname{Sym}^{p}(\mathcal{V})\defeq T_p/R_{\vee}$.
\item Set $A_p\defeq\operatorname{span}\{\lambda_{\sigma}\mathbf{z}-\sign(\sigma)\mathbf{z}:\sigma\in G_p,\mathbf{z}\in T_p\}$ and $D_p\defeq\operatorname{span}\{\mathbf{v}_1\otimes\cdots\otimes\mathbf{v}_p:\exists i<j,\mathbf{v}_i=\mathbf{v}_j\}$.\\
The \hlt{exterior power} is $\bigwedge^{p}\mathcal{V}\defeq T_p/(A_p+D_p)$.
\item For quotient map $q_{\vee}$, the \hlt{symmetric product} is $q_{\vee}(\mathbf{v}_1\otimes\cdots\otimes\mathbf{v}_p)=\mathbf{v}_1\vee\cdots\vee\mathbf{v}_p$
\item For quotient map $q_{\wedge}$, the \hlt{exterior product} is $q_{\wedge}(\mathbf{v}_1\otimes\cdots\otimes\mathbf{v}_p)=\mathbf{v}_1\wedge\cdots\wedge\mathbf{v}_p$.
\end{enumerate}
\end{definition}

\begin{theorem} \label{thm:symmetric-exterior-universal-properties} \hlt{Universal Properties of Symmetric and Exterior Powers}\\
Let $\mathcal{U}$ be a vector space, $M=\operatorname{Mult}_{\mathbb{F}}(\mathcal{V},\ldots,\mathcal{V};\mathcal{U})$, and $f^\sigma(\mathbf{v}_1,\ldots,\mathbf{v}_p)\defeq f(\mathbf{v}_{\sigma(1)},\ldots,\mathbf{v}_{\sigma(p)})$.
\begin{enumerate}[label=(\roman*)]
\setlength{\itemsep}{0pt}
\item $\LL(\operatorname{Sym}^{p}(\mathcal{V}),\mathcal{U})\cong M_{\vee}$, where $M_{\vee}\defeq\{f\in M:\forall\sigma\in\mathfrak{S}_p,\ f^\sigma=f\}$.
\item $\LL(\bigwedge^{p}\mathcal{V},\mathcal{U})\cong M_{\wedge}$, where $M_{\wedge}\defeq\{f\in M:\forall\sigma\in\mathfrak{S}_p,\ f^\sigma=\sign(\sigma)f,\ \forall(\mathbf{v}_1,\ldots,\mathbf{v}_p)$,\\
$(\exists i<j,\mathbf{v}_i=\mathbf{v}_j)\Rightarrow f(\mathbf{v}_1,\ldots,\mathbf{v}_p)=0\}$.
\end{enumerate}
\begin{proof}
\begin{enumerate}[label=(\roman*)]
\setlength{\itemsep}{0pt}
\item
Let $t:\mathcal{V}^{\times p}\to\mathcal{V}^{\otimes p}$ and $q_{\vee}:\mathcal{V}^{\otimes p}\to\operatorname{Sym}^{p}(\mathcal{V})$ be canonical.\\
The map $\Phi_{\vee}(T)=Tq_{\vee}t$ sends $\LL(\operatorname{Sym}^{p}(\mathcal{V}),\mathcal{U})$ into $M_{\vee}$.\\
Conversely, for $f\in M_{\vee}$, there is a unique $\widetilde{f}:\mathcal{V}^{\otimes p}\to\mathcal{U}$ with $f=\widetilde{f}t$; since $\widetilde{f}(R_{\vee})=0$, it factors uniquely through $q_{\vee}$ as $\widetilde{f}=\overline{f}q_{\vee}$, where $\overline{f}:\operatorname{Sym}^{p}(\mathcal{V})\to\mathcal{U}$. Thus $f\mapsto\overline{f}$ is $\Phi_{\vee}^{-1}$.
\item Let $t:\mathcal{V}^{\times p}\to\mathcal{V}^{\otimes p}$ and $q_{\wedge}:\mathcal{V}^{\otimes p}\to\bigwedge^{p}\mathcal{V}$ be canonical.\\
The map $\Phi_{\wedge}(T)=Tq_{\wedge}t$ sends $\LL(\bigwedge^{p}\mathcal{V},\mathcal{U})$ into $M_{\wedge}$. Conversely, for $f\in M_{\wedge}$, the induced $\widetilde{f}:\mathcal{V}^{\otimes p}\to\mathcal{U}$ satisfies $\widetilde{f}(R_{\wedge})=0$, hence factors uniquely as $\widetilde{f}=\overline{f}q_{\wedge}$; thus $f\mapsto\overline{f}$ is $\Phi_{\wedge}^{-1}$.
\end{enumerate}
\hlt{Structural Commentary}: Symmetric and exterior powers are tensor quotients by permutation relations.\\
\hlt{Key Insight}: Symmetric and alternating multilinear maps are ordinary linear maps out of the right quotient.
\end{proof}
\end{theorem}

\begin{theorem} \label{thm:symmetric-exterior-dimensions} \hlt{Dimensions of Symmetric and Exterior Powers}\\
Let $\dim\mathcal{V}=n<\infty$ and $p\in\mathbb{Z}_{\geq0}$. Set $I_{\vee}=\{1\leq i_1\leq\cdots\leq i_p\leq n\}$ and $I_{\wedge}=\{1\leq i_1<\cdots<i_p\leq n\}$.
\begin{enumerate}[label=(\roman*)]
\setlength{\itemsep}{0pt}
\item $\dim\operatorname{Sym}^{p}(\mathcal{V})=\#I_{\vee}=\binom{n+p-1}{p}$.
\item $\dim\bigwedge^{p}\mathcal{V}=\#I_{\wedge}=\binom{n}{p}$, with $\binom{n}{p}=0$ for $p>n$.
\end{enumerate}
\begin{proof}
Let $E=(\mathbf{e}_1,\ldots,\mathbf{e}_n)$ be a basis. For $I=(i_1,\ldots,i_p)$, set $\mathbf{e}_{I}^{\vee}=\mathbf{e}_{i_1}\vee\cdots\vee\mathbf{e}_{i_p}$ and $\mathbf{e}_{I}^{\wedge}=\mathbf{e}_{i_1}\wedge\cdots\wedge\mathbf{e}_{i_p}$.
\begin{enumerate}[label=(\roman*)]
\setlength{\itemsep}{0pt}
\item
$R_{\vee}$ identifies each word with its nondecreasing representative. Thus $B_{\vee}=\{\mathbf{e}_{I}^{\vee}:I\in I_{\vee}\}$ is the quotient basis, and $\dim\operatorname{Sym}^{p}(\mathcal{V})=\#B_{\vee}=\#I_{\vee}=\binom{n+p-1}{p}$.
\item
$R_{\wedge}$ kills repeated words and identifies each remaining word with its increasing representative up to sign. Thus $B_{\wedge}=\{\mathbf{e}_{I}^{\wedge}:I\in I_{\wedge}\}$ is the quotient basis, and $\dim\bigwedge^{p}\mathcal{V}=\#B_{\wedge}=\#I_{\wedge}=\binom{n}{p}$.
\end{enumerate}
\hlt{Structural Commentary}: Symmetric powers count repetitions; exterior powers forbid repetitions.\\
\hlt{Key Insight}: Multisets give symmetric dimensions, subsets give exterior dimensions.
\end{proof}
\end{theorem}

\begin{proposition} \label{prop:symmetrization-map} \hlt{Symmetrization in Characteristic Zero}\\
Assume $\operatorname{char}(\mathbb{F})=0$. Define $S,A\in\LL(\mathcal{V}^{\otimes p})$ by $S\defeq\frac{1}{p!}\sum\limits_{\sigma\in\mathfrak{S}_p}\lambda_{\sigma}$ and $A\defeq\frac{1}{p!}\sum\limits_{\sigma\in\mathfrak{S}_p}\sign(\sigma)\lambda_{\sigma}$.
\begin{enumerate}[label=(\roman*)]
\setlength{\itemsep}{0pt}
\item $S^2=S$ and $\operatorname{im}S=\operatorname{ST}^{p}(\mathcal{V})$.
\item $A^2=A$ and $\operatorname{im}A=\operatorname{AT}^{p}(\mathcal{V})$.
\end{enumerate}
\begin{proof}
\begin{enumerate}[label=(\roman*)]
\setlength{\itemsep}{0pt}
\item For $\tau\in\mathfrak{S}_p$, $\lambda_{\tau}S=p!^{-1}\sum_{\sigma}\lambda_{\tau\sigma}=S$, hence $\operatorname{im}S\subseteq\operatorname{ST}^{p}(\mathcal{V})$.\\
If $\mathbf{z}\in\operatorname{ST}^{p}(\mathcal{V})$, then $S\mathbf{z}=\mathbf{z}$, so $\operatorname{im}S=\operatorname{ST}^{p}(\mathcal{V})$ and $S^2=S$.
\item For $\tau\in\mathfrak{S}_p$, $\lambda_{\tau}A=p!^{-1}\sum_{\sigma}\sign(\sigma)\lambda_{\tau\sigma}=\sign(\tau)A$, hence $\operatorname{im}A\subseteq\operatorname{AT}^{p}(\mathcal{V})$.\\
If $\mathbf{z}\in\operatorname{AT}^{p}(\mathcal{V})$, then $A\mathbf{z}=p!^{-1}\sum_{\sigma}\sign(\sigma)^2\mathbf{z}=\mathbf{z}$, so $\operatorname{im}A=\operatorname{AT}^{p}(\mathcal{V})$ and $A^2=A$.
\end{enumerate}
\hlt{Structural Commentary}: Averaging gives symmetry projections.\\
\hlt{Key Insight}: The factor $1/p!$ makes the averaging map fix the tensors it projects onto.
\end{proof}
\end{proposition}

\begin{definition} \label{def:graded-tensor-symmetric-exterior-algebras} \hlt{Tensor, Symmetric, and Exterior Algebras}\\
Let $\dim\mathcal{V}<\infty$.
\begin{enumerate}[label=(\roman*)]
\setlength{\itemsep}{0pt}
\item The \hlt{tensor algebra} is $\mathcal{T}(\mathcal{V})\defeq\bigoplus_{p\geq0}\mathcal{V}^{\otimes p}$, with $\mathcal{V}^{\otimes0}\defeq\mathbb{F}$ and product induced by tensor concatenation.
\item The \hlt{symmetric algebra} is $\operatorname{Sym}(\mathcal{V})\defeq\bigoplus_{p\geq0}\operatorname{Sym}^{p}(\mathcal{V})$, with product induced by $\vee$.
\item The \hlt{exterior algebra} is $\bigwedge\mathcal{V}\defeq\bigoplus_{0\leq p\leq\dim\mathcal{V}}\bigwedge^{p}\mathcal{V}$, with product induced by $\wedge$.
\item Each is \hlt{graded}: degree $p$ times degree $q$ lies in degree $p+q$.
\end{enumerate}
\end{definition}
